\documentclass[a4paper,12pt,twoside]{memoir}

% English
\usepackage[english]{babel}
\selectlanguage{english}
\usepackage[utf8]{inputenc}
\usepackage[T1]{fontenc}
\usepackage{lmodern} % Scalable font
\usepackage{microtype}
\usepackage{placeins}
\usepackage{listings}

\RequirePackage{booktabs}
\RequirePackage[table]{xcolor}
\RequirePackage{xtab}
\RequirePackage{multirow}
\PassOptionsToPackage{hyphens}{url}

% Links
\usepackage[colorlinks]{hyperref}
\hypersetup{
	allcolors = {red}
}

% Ecuaciones
\usepackage{amsmath}

% Rutas de fichero / paquete
\newcommand{\ruta}[1]{{\sffamily #1}}

% Párrafos
\nonzeroparskip

% Huérfanas y viudas
\widowpenalty100000
\clubpenalty100000

% Evitar solapes en el header
\nouppercaseheads

% Imagenes
\usepackage{graphicx}
\newcommand{\imagen}[2]{
	\begin{figure}[!h]
		\centering
		\includegraphics[width=0.9\textwidth]{#1}
		\caption{#2}\label{fig:#1}
	\end{figure}
	\FloatBarrier
}

\newcommand{\imagenflotante}[2]{
	\begin{figure}%[!h]
		\centering
		\includegraphics[width=0.9\textwidth]{#1}
		\caption{#2}\label{fig:#1}
	\end{figure}
}



% El comando \figura nos permite insertar figuras comodamente, y utilizando
% siempre el mismo formato. Los parametros son:
% 1 -> Porcentaje del ancho de página que ocupará la figura (de 0 a 1)
% 2 --> Fichero de la imagen
% 3 --> Texto a pie de imagen
% 4 --> Etiqueta (label) para referencias
% 5 --> Opciones que queramos pasarle al \includegraphics
% 6 --> Opciones de posicionamiento a pasarle a \begin{figure}
\newcommand{\figuraConPosicion}[6]{%
  \setlength{\anchoFloat}{#1\textwidth}%
  \addtolength{\anchoFloat}{-4\fboxsep}%
  \setlength{\anchoFigura}{\anchoFloat}%
  \begin{figure}[#6]
    \begin{center}%
      \Ovalbox{%
        \begin{minipage}{\anchoFloat}%
          \begin{center}%
            \includegraphics[width=\anchoFigura,#5]{#2}%
            \caption{#3}%
            \label{#4}%
          \end{center}%
        \end{minipage}
      }%
    \end{center}%
  \end{figure}%
}

%
% Comando para incluir imágenes en formato apaisado (sin marco).
\newcommand{\figuraApaisadaSinMarco}[5]{%
  \begin{figure}%
    \begin{center}%
    \includegraphics[angle=90,height=#1\textheight,#5]{#2}%
    \caption{#3}%
    \label{#4}%
    \end{center}%
  \end{figure}%
}

% Comando para incluir imágenes en formato normal (sin marco).
\newcommand{\figuraNormalSinMarco}[5]{%
  \begin{figure}%
    \begin{center}%
    \includegraphics[angle=0,height=#1\textheight,#5]{#2}%
    \caption{#3}%
    \label{#4}%
    \end{center}%
  \end{figure}%
}

% Para las tablas
\newcommand{\otoprule}{\midrule [\heavyrulewidth]}
%
% Nuevo comando para tablas pequeñas (menos de una página).
\newcommand{\tablaSmall}[5]{%
 \begin{table}
  \begin{center}
   \rowcolors {2}{gray!35}{}
   \begin{tabular}{#2}
    \toprule
    #4
    \otoprule
    #5
    \bottomrule
   \end{tabular}
   \caption{#1}
   \label{tabla:#3}
  \end{center}
 \end{table}
}

%
% Nuevo comando para tablas pequeñas (menos de una página).
\newcommand{\tablaSmallSinColores}[5]{%
 \begin{table}[H]
  \begin{center}
   \begin{tabular}{#2}
    \toprule
    #4
    \otoprule
    #5
    \bottomrule
   \end{tabular}
   \caption{#1}
   \label{tabla:#3}
  \end{center}
 \end{table}
}

\newcommand{\tablaApaisadaSmall}[5]{%
\begin{landscape}
  \begin{table}
   \begin{center}
    \rowcolors {2}{gray!35}{}
    \begin{tabular}{#2}
     \toprule
     #4
     \otoprule
     #5
     \bottomrule
    \end{tabular}
    \caption{#1}
    \label{tabla:#3}
   \end{center}
  \end{table}
\end{landscape}
}

%
% Nuevo comando para tablas grandes con cabecera y filas alternas coloreadas en gris.
\newcommand{\tabla}[6]{%
  \begin{center}
    \tablefirsthead{
      \toprule
      #5
      \otoprule
    }
    \tablehead{
      \multicolumn{#3}{l}{\small\sl continúa desde la página anterior}\\
      \toprule
      #5
      \otoprule
    }
    \tabletail{
      \hline
      \multicolumn{#3}{r}{\small\sl continúa en la página siguiente}\\
    }
    \tablelasttail{
      \hline
    }
    \bottomcaption{#1}
    \rowcolors {2}{gray!35}{}
    \begin{xtabular}{#2}
      #6
      \bottomrule
    \end{xtabular}
    \label{tabla:#4}
  \end{center}
}

%
% Nuevo comando para tablas grandes con cabecera.
\newcommand{\tablaSinColores}[6]{%
  \begin{center}
    \tablefirsthead{
      \toprule
      #5
      \otoprule
    }
    \tablehead{
      \multicolumn{#3}{l}{\small\sl continúa desde la página anterior}\\
      \toprule
      #5
      \otoprule
    }
    \tabletail{
      \hline
      \multicolumn{#3}{r}{\small\sl continúa en la página siguiente}\\
    }
    \tablelasttail{
      \hline
    }
    \bottomcaption{#1}
    \begin{xtabular}{#2}
      #6
      \bottomrule
    \end{xtabular}
    \label{tabla:#4}
  \end{center}
}

%
% Nuevo comando para tablas grandes sin cabecera.
\newcommand{\tablaSinCabecera}[5]{%
  \begin{center}
    \tablefirsthead{
      \toprule
    }
    \tablehead{
      \multicolumn{#3}{l}{\small\sl continúa desde la página anterior}\\
      \hline
    }
    \tabletail{
      \hline
      \multicolumn{#3}{r}{\small\sl continúa en la página siguiente}\\
    }
    \tablelasttail{
      \hline
    }
    \bottomcaption{#1}
  \begin{xtabular}{#2}
    #5
   \bottomrule
  \end{xtabular}
  \label{tabla:#4}
  \end{center}
}



\definecolor{cgoLight}{HTML}{EEEEEE}
\definecolor{cgoExtralight}{HTML}{FFFFFF}

%
% Comando para el texto en formato JSON
% Source: https://tex.stackexchange.com/questions/83085/how-to-improve-listings-display-of-json-files

% Switch used as state variable
\makeatletter
\newif\ifisvalue@json

\newcommand\jsonkey{\color{purple}}
\newcommand\jsonvalue{\color{cyan}}
\newcommand\jsonnumber{\color{orange}}
\definecolor{background}{HTML}{EEEEEE}

\lstdefinelanguage{json}{
    tabsize=4,
    showstringspaces=false,
    breaklines=true,
    keywords={false,true},
    alsoletter=0123456789.,
    morestring=[s]{"}{"},
    stringstyle=\jsonkey\ifisvalue@json\jsonvalue\fi,
    MoreSelectCharTable=\lst@DefSaveDef{`:}\colon@json{\enterMode@json},
    MoreSelectCharTable=\lst@DefSaveDef{`,}\comma@json{\exitMode@json{\comma@json}},
    MoreSelectCharTable=\lst@DefSaveDef{`\{}\bracket@json{\exitMode@json{\bracket@json}},
    basicstyle=\scriptsize\ttfamily,
    backgroundcolor=\color{background}
}

% Enter "value" mode after encountering a colon
\newcommand\enterMode@json{%
    \colon@json%
    \ifnum\lst@mode=\lst@Pmode%
        \global\isvalue@jsontrue%
    \fi
}

% Leave "value" mode: either we hit a comma, or the value is a nested object
\newcommand\exitMode@json[1]{#1\global\isvalue@jsonfalse}

\lst@AddToHook{Output}{%
    \ifisvalue@json%
        \ifnum\lst@mode=\lst@Pmode%
            \def\lst@thestyle{\jsonnumber}%
        \fi
    \fi
    % Override by keyword style if a keyword is detected
    \lsthk@DetectKeywords% 
}

\makeatother

%
% Nuevo comando para tablas grandes sin cabecera.
\newcommand{\tablaSinCabeceraConBandas}[5]{%
  \begin{center}
    \tablefirsthead{
      \toprule
    }
    \tablehead{
      \multicolumn{#3}{l}{\small\sl continúa desde la página anterior}\\
      \hline
    }
    \tabletail{
      \hline
      \multicolumn{#3}{r}{\small\sl continúa en la página siguiente}\\
    }
    \tablelasttail{
      \hline
    }
    \bottomcaption{#1}
    \rowcolors[]{1}{cgoExtralight}{cgoLight}

  \begin{xtabular}{#2}
    #5
   \bottomrule
  \end{xtabular}
  \label{tabla:#4}
  \end{center}
}


\graphicspath{ {./img/} }

% Capítulos
\chapterstyle{bianchi}
\newcommand{\capitulo}[2]{
	\setcounter{chapter}{#1}
	\setcounter{section}{0}
	\chapter*{#2}
	\addcontentsline{toc}{chapter}{#1. #2}
	\markboth{#2}{#2}
}

% Apéndices
\renewcommand{\appendixname}{Appendix}
\renewcommand*\cftappendixname{\appendixname}

\newcommand{\apendice}[1]{
	%\renewcommand{\thechapter}{A}
	\chapter{#1}
}

\renewcommand*\cftappendixname{\appendixname\ }

% Formato de portada
\makeatletter
\usepackage{xcolor}
\newcommand{\tutor}[1]{\def\@tutor{#1}}
\newcommand{\course}[1]{\def\@course{#1}}
\definecolor{cpardoBox}{HTML}{E6E6FF}
\def\maketitle{
  \null
  \thispagestyle{empty}
  % Cabecera ----------------
\begin{center}%
	{\noindent\Huge Universities of Burgos, León and Valladolid}\vspace{.5cm}%
	
	{\noindent\Large Master's degree}\vspace{.5cm}%
	
	{\noindent\Huge \textbf{Business Intelligence and Big Data in Cyber-Secure Environments}}\vspace{.5cm}%
\end{center}%

\begin{center}%
	\includegraphics[height=3cm]{img/escudoUBU} \hspace{1cm}
	\includegraphics[height=3cm]{img/escudoUVA} \hspace{1cm}
	\includegraphics[height=3cm]{img/escudoULE} \vspace{0.1cm}%
\end{center}%

  \vfill
  % Título proyecto y escudo informática ----------------
  \colorbox{cpardoBox}{%
    \begin{minipage}{0.9\textwidth}
      \vspace{.5cm}\large
      \begin{center}
      \textbf{Thesis of the Master's degree in Business Intelligence and Big Data in Cyber-Secure Environments}\vspace{.5cm}\\
      \textbf{\Large\@title{}}
      \end{center}
      \vspace{.2cm}
    \end{minipage}

  }%
  \hfill
  \vfill
  % Datos de alumno, curso y tutores ------------------
  \begin{center}%
  {%
    \noindent\Large
    Presented by \@author{}\\ 
    in University of Burgos --- \@date{}\\
    Supervisor: \@tutor{}\\
  }%
  \end{center}%
  \null
  \cleardoublepage
  }
\makeatother

\newcommand{\nombre}{Adrián Riesco Valbuena} %%% cambio de comando

% Datos de portada
\title{Extraction, transformation, loading and visualization of combined Twitter and Spotify data in a scalable architecture}
\author{\nombre}
\tutor{Álvar Arnaiz González}
\date{\today}

\begin{document}

\maketitle


\newpage\null\thispagestyle{empty}\newpage


%%%%%%%%%%%%%%%%%%%%%%%%%%%%%%%%%%%%%%%%%%%%%%%%%%%%%%%%%%%%%%%%%%%%%%%%%%%%%%%%%%%%%%%%
\thispagestyle{empty}


\noindent
\begin{center}%
	{\noindent\Huge Universities of Burgos, León and Valladolid}\vspace{.5cm}%
	
\begin{center}%
	\includegraphics[height=3cm]{img/escudoUBU} \hspace{1cm}
	\includegraphics[height=3cm]{img/escudoUVA} \hspace{1cm}
	\includegraphics[height=3cm]{img/escudoULE} \vspace{1cm}%
\end{center}%

	{\noindent\Large \textbf{Master's degree in Business Intelligence and Big Data in Cyber-Secure Environments}}\vspace{.5cm}%
\end{center}%



\noindent D. Álvar Arnaiz González, professor of the department named Computer Engineering, area named Computer Languages and Systems.

\noindent Exposes:

\noindent That the student Mr. \nombre, with DNI 71462231N, has completed the Thesis of the Master in Business Intelligence and Big Data in Cyber-Secure Environments titled ``Extraction, transformation, loading and visualization of combined Twitter and Spotify data in a scalable architecture''. 

\noindent And that thesis has been carried out by the student under the direction of the undersigned, by virtue of which its presentation and defense is authorized.
\begin{center} %\large
In Burgos, {\today}
\end{center}

\vfill\vfill\vfill

% Tutor
\begin{minipage}{0.45\textwidth}
\begin{flushleft} %\large
Approval of the Supervisor:\\[2cm]
D. Álvar Arnaiz González
\end{flushleft}
\end{minipage}


\newpage\null\thispagestyle{empty}\newpage

\frontmatter

% Abstract en castellano
\renewcommand*\abstractname{Resumen}
\begin{abstract}
*****¿ELIMINAR ABSTRACT EN ESPAÑOL?*****
\end{abstract}

\renewcommand*\abstractname{Descriptores}
\begin{abstract}
Palabras separadas por comas que identifiquen el contenido del proyecto Ej: servidor web, buscador de vuelos, android \ldots
\end{abstract}

\clearpage

% Abstract en inglés
\renewcommand*\abstractname{Abstract}
\begin{abstract}
This project consists of the development of an \textbf{Extraction, Transformation and Loading} process to capture data from \textbf{Twitter} and \textbf{Spotify} APIs using software tools from the Big Data field, ensuring its scalability. Data from both sources ha been combined, processed and stored in a Data Warehouse to be served to the user in a customized front-end. All services have been encapsulated in containers within the same network.
\end{abstract}

\renewcommand*\abstractname{Keywords}
\begin{abstract}
Airflow, Apache, BigData, Bootstrap, Cassandra, Chart.js, Datatables, Docker, Docker Compose, ETL, Flask, Spark, Spotify, Twitter.
\end{abstract}

\clearpage

% Indices
\tableofcontents

\clearpage

\listoffigures

\clearpage

\listoftables
\clearpage

\mainmatter

\addcontentsline{toc}{part}{Memory}
\part*{Memory}

\capitulo{1}{Introduction}

% Description of the work, the structure of the memory and the rest of the material delivered.

\nonzeroparskip Social networks are currently a fundamental aspect of society. People usually use social networks to share experiences, opinions, and aspects of their lives and interact with other people. Using social networks as a data source we can access a huge amount of information and be able to build accurate analyses on practically any topic.

\nonzeroparskip On the other hand, another aspect that has gradually permeated our society is the concept of subscriptions to services, be it music, movies, games, or almost any concept that we can think of. Not so many years ago, the concept of paying for subscriptions to services, where you do not actually own the content you pay for and instead get temporary access whose duration is defined by how long you continue to pay for the subscription, was relegated to very specific services and was not nearly as globalized as it is today.

\nonzeroparskip The global acceptance of subscription as a service is reflected in social networks, where users can comment on the different music, movie, and game platforms, turning each new release into a social phenomenon. In order to take advantage of both worlds, this project uses the social network Twitter to obtain information about the latest music listened to by users (by searching for a particular hashtag) and then consult the data of the song and the artist involved that Spotify, a platform based on music as a service, has.

\nonzeroparskip For the development of the project, a series of scripts have been designed for the execution of the ETL process (Extract, Transform and Load) aimed at consuming the data from both APIs, processing them and loading them in the data warehouse, for their subsequent sending to a custom front end, all based on recognized tools in the field of Big Data, such as Apache Airflow, Apache Spark, Apache Cassandra, Docker Compose and Flask, among others. The project has been organized using an agile methodology with two-week sprints.
\capitulo{2}{Project objectives}

\nonzeroparskip The initial objectives through which the use case was built consisted in the following points:
\begin{itemize}
	\item Ability to obtain data in real time.
	\item Combination of at least two different data sources.
	\item Potential to scale in both technology and data volume.
	\item Implication of recognized technologies within the field of Big Data.
	\item Use of open source tools to avoid potential costs.
\end{itemize}

\nonzeroparskip After a research, the author designed the use case and built the objectives of the projects:
\begin{itemize}
	\item Build a pipeline to gather information about last songs listened from the Twitter API. The name of the endpoint queried is \textit{recent search}.
	\item Find information about the songs (name, artist and audio features) through the Spotify API. The names of the endpoints queried are \textit{search for item} and \textit{get tracks' audio features}.
	\item Execute all the ETL\footnote{ETL is the acronym for Extract, Transform and Load, the three phases for data processing} process in \textit{Apache Spark}.
	\item Store all the data in a data warehouse under a known technology, \textit{Apache Cassandra}.
	\item Visualize the information in a custom front-end and back-end created with \textit{Bootstrap} and \textit{Flask}.
	\item Orchestrate all the data flow with \textit{Apache Airflow}.
	\item Develop the project with \textit{Docker} and \textit{Docker Compose} to ensure deployment through heterogeneous environments.
\end{itemize}

\nonzeroparskip Through these global objectives, the low-level functional and technical requirements were specified. More information about these requirements is included in the appendix.
\capitulo{3}{Theoretical concepts}

%Theoretical concepts of \LaTeX \footnote{Example of footnote}.
%\subsection{Subsection}
%\subsubsection{Subsubsection}
%Use of cite \cite{wiki:latex}, \cite{koza92}.
%\imagen{img/escudoInfor}{Image caption}
%Lists
%\begin{itemize}
%Enumerate.
%\begin{enumerate}
%Description.
%\begin{description}

\nonzeroparskip In this section are covered the theoretical concepts in which the project has been based. All concepts are described in a detailed an simple way since this master's degree can be aimed at technical and non-technical students.

\section{API}

\nonzeroparskip Section explaining API concepts.

\section{Orchestrator}

\nonzeroparskip Section explaining Flow Orchestrator -> Airflow.

\section{NoSQL Databases}

\nonzeroparskip Section explaining NoSQL Databases.

\section{Containers}

\nonzeroparskip Section explaining Containers.

\section{Continuous Integration / Continuous Delivery}

\nonzeroparskip Section explaining CI/CD.

\section{Template engines}

\nonzeroparskip Section explaining Template engines -> Jinja.

\section{Web Server Gateway Interface}

\nonzeroparskip Section explaining Web Server Gateway Interface (WSGI).


\section{Tables}

\nonzeroparskip TablaSmall.

\tablaSmall{Tools and technologies used}{l c c c c}{herramientasportipodeuso}
{ \multicolumn{1}{l}{Tools} & App AngularJS & API REST & BD & Memoria \\}{ 
HTML5 & X & & &\\
CSS3 & X & & &\\
BOOTSTRAP & X & & &\\
\TeX{}Maker & & & & X\\
Astah & & & & X\\
} 

\capitulo{4}{Techniques and tools}

\nonzeroparskip In this section are presented the methodological techniques and development tools used to carry out the project.

\section{GitHub}

\nonzeroparskip GitHub is the repository where the project was uploaded and its evolution was tracked.

\section{APIs}

\nonzeroparskip During this project there were used APIs from two different providers to gather the information: Twitter API and Spotify API.

\section{Postman}

\nonzeroparskip Postman is a tool that allows the user to build and use APIs in a simple way. Some of its characteristics are the API repository (easier storage, cataloging and collaboration), the availability of tools to help in the API design, testing and documentation, the workspaces to organize the work, and built-in integrations with tools such as GitHub, Azure DevOps, Jenkins, Splunk, Slack and Microsoft Teams. In addition, Postman is based on open source technologies, which provides the ability to be easily extended~\cite{postman}.

\section{Apache Airflow}

\nonzeroparskip Apache Airflow is a service orchestrator that allows you to plan, manage, and monitor workflows~\cite{airflow}. It was created in 2014 by Airbnb with the aim of handling the company's huge data flows, and published in 2015 under an open source license. In March 2016 the project joined the Apache Software Foundations incubator and was published as a top level project in 2019.

\nonzeroparskip Airflow is used to automate jobs by breaking them down into smaller tasks. For example, this project uses Airflow to automate the ETL process that consumes data from Twitter and Spotify, processes it, and serves it to the user. Among the main features of Airflow are scalability and ease of integration with other tools.

\nonzeroparskip The main element used by Airflow are the Directed Acyclic Graphs or DAGs, which are groups of tasks connected to each other through dependencies like the nodes of a graph. The word \textit{direct} indicates that the existing relationships in the graph must only have one direction (bidirectional relationships between nodes or tasks are not allowed), while the word \textit{acyclic} means that cycles cannot exist in the graph (nodes or tasks cannot be executed more than once). Tasks are defined by means of an operator and there is a very extensive library with operators that allow defining a wide variety of services such as BashOperator, to execute Bash commands, or SparkSubmitOperator, to submit a task to Spark. Regarding the programming language, Python is the one in which DAGs are developed.

\nonzeroparskip Airflow allows you to have control of the tasks executed through a record of their executions, the time, current or final status and the generated logs. In addition, it allows certain parameters to be associated with each task, such as, for example, the maximum execution time allowed.

\section{Apache Spark}

\nonzeroparskip Apache Spark is...

\section{Cassandra}

\nonzeroparskip Cassandra is a NoSQL database that...

\section{Flask}

\nonzeroparskip Flask is...

\section{Bootstrap}

\nonzeroparskip Bootstrap is...

\section{Docker}

\nonzeroparskip Docker is...

\section{Docker Compose}

\nonzeroparskip Docker Compose is...

\capitulo{5}{Relevant aspects of the project}

%This section aims to collect the most interesting aspects of the development of the project, commented by its authors.
%It must include from the exposition of the life cycle used, to the most relevant details of the analysis, design and implementation phases.
%It is sought that it is not a mere operation of copying and pasting diagrams and extracts from the source code, but that the solution paths that have been taken are really justified, especially those that are not trivial.
%It may be the most appropriate place to document the most interesting aspects of the design and implementation, with a greater emphasis on aspects such as the type of architecture chosen, the indexes of the database tables, normalization and denormalization, distribution in files3, business rules within databases (EDVHV GH GDWRV DFWLYDV), development aspects related to the WWW...
%This section must become the summary of the practical experience of the project, and by itself justifies that the report becomes a useful document, a reference source for authors, tutors and future students.

\nonzeroparskip The first step of the project was the feasibility and viability analysis of the concept devised. The author was looking to use two data sources with:

\begin{itemize}
	\item Real and updated data, preferable related to the social interest.
	\item The possibility of getting a stream data flow.
	\item The potential to combine both to get an added value.
\end{itemize}

\nonzeroparskip Considering the previous points, the author found an interesting option on Twitter and Spotify providers. Both of them provides solid APIs for a fluid development and have the characteristics needed to combine the data collected. Consequently, the author designed the following use case:

\begin{enumerate}
	\item The Twitter API is consulted to gather the \textit{tweets} with the hashtag \textit{\#NowPlaying}.
	\item The tweet is cleaned, removing the stopwords and the other hashtags and getting the song name and artist as isolated as possible.
	\item The Spotify API is consulted to gather the information of the song identified.
	\item The vector values of the cleaned Twitter data and the name of the song returned by Spotify are compared to ensure they are the same.
	\item The data is moved to the database, ready to be stored and visualized.
\end{enumerate}

\nonzeroparskip During the design phase, the author analyzed the output of both APIs using Postman.


\nonzeroparskip The project development was undertaken following an Agile methodology.\\
\capitulo{6}{Related works}


\nonzeroparskip Both the software tools used in the project, widely recognized within the Big Data industry, and the data sources, from large cap and renowned companies, have been applied to countless projects throughout their lifetime.

\nonzeroparskip During the planning and study phase of the APIs and tools, the author discovered several projects that have combined both elements in different ways. Although no project has been found that combines both APIs and involves a data flow similar to that of the present project, the author has decided to mention those considered most interesting, either because of their data flow, their usefulness or the originality of the idea:

\begin{itemize}
	\item \textbf{Tweetpy}~\footnote{\url{https://www.tweepy.org/}}. Created in 2009 by Joshua Roesslein, this project contributes a Python library to interact with the Twitter API, simplifying interactions during development. Despite considering it as an alternative, the author decided not to make use of it in order to learn directly how to interact with the API. However, in the case of having to interact with user accesses, it would probably have been chosen.
	\item \textbf{Spotipy}~\footnote{\url{https://spotipy.readthedocs.io/}}. Created in 2014 by Paul Lamere, this project provides a Python library to interact with the Spotify API. As with the previous library, it was decided not to use it to learn from direct interaction with the API, since the data flow of the project does not involve complex access permissions.
	\item \textbf{SpotaTweet}~\footnote{\url{https://github.com/twitterdev/spotatweet}}. Although its data flow is not analogous to that of the project, its use case bears some similarity, with a collection of tweets with the hashtag \texttt{\#NowPlaying} and a subsequent search for the song listened to on Spotify with the goal of displaying a preview to the user.
	\item \textbf{Divide for Spotify}~\footnote{\url{https://divideforspotify.com/}}. This project makes use of the Spotify API to divide the songs that have been liked among the playlists of the user in question. This project has been chosen for its originality, as it takes advantage of the API to add a very interesting functionality to the platform, automatically dividing the songs of a playlist to choose from a selection of playlists chosen by the user.
	\item \textbf{Spotify to Twitter}~\footnote{\url{https://github.com/transitive-bullshit/spotify-to-twitter}}. This project combines both APIs to automate a Twitter account to publish tweets with songs from a Spotify playlist. The project is based on \texttt{Node.js}, an open source, cross-platform execution environment for the JavaScript-based server layer.
\end{itemize}

\nonzeroparskip The projects mentioned above are only a fragment of those encountered during the initial analysis. Due to the breadth of Big Data domains, the use cases are innumerable and specialized tools can be of use in a number of different environments and industries.

\nonzeroparskip There are also a large number of scientific articles dedicated to the study of data from both social networks, especially Twitter since it is a social network with a microblogging format that allows for interesting analysis. For example, one of the most well-known projects is the sentiment analysis of tweets, which is often combined with data mining on a given topic.

\nonzeroparskip In the same way, all of the software tools used have a wealth of resources, strong communities and scientific articles where they are applied within their scope of operation. The author decided not to include in this section the scientific articles found during the study and analysis phase because their use cases are not particularly relevant to the project.
\capitulo{7}{Conclusions and future work lines}

% Every project must include the conclusions derived from its development. These can be of a different nature, depending on the type of project, but normally there will be a set of conclusions related to the results of the project and a set of technical conclusions.
% In addition, it is very useful to make a critical report indicating how the project can be improved, or how work can continue along the lines of the completed project.

% Next steps: clean data, check songs, ask for last id tweet instead of get last n (if no rate limits), add group by, add global metrics, add range of metrics, group by artist, optimize the process (duration), button to load from historic...
% Create MV, developer keys, change GitHub name, change project name, video demo, add summary, descriptors, add figures.

% Alvar:
% - ¿Eliminar resumen y descriptores en español?
% - ¿Título del TFM y del GitHub?
% Twitter and Spotify data extraction, transformation, loading and visualization using Big Data and development tools.
% - ¿Claves de desarrollador para la MV?

% - Más problemas? Comprensión inicial de Docker? Configuración inicial de Airflow?
% - Capturas de UIs -> Explicación breve de los elementos de cada vista.
% - Captura de Aiflow
% - Agradecimientos.
% - ¿?

\nonzeroparskip Given the requirements established in the planning phase, the final result is considered a success. The author has learned how to handle the tools included in the scope and how to interconnect them to achieve a complete ETL flow with a simple but effective visualization layer.

\nonzeroparskip In case of continuing with the presented work, the author has performed an exercise of analysis and identification of the next steps and improvements that could be implemented:
\begin{itemize}
	\item \textbf{Improve the percentage of correctly identified tracks.} This action could be addressed in several ways, including: discarding tweets with characters outside the Latin alphabet (not only specific characters), removing stop words in several languages so as not to have to discard all characters not included in the Latin alphabet, or increasing the number of tracks returned by Spotify and performing a comparative method (bag of words, TI-IDF, etc.) between their titles and the full text of the tweet. This would improve the reliability of the database and give more confidence when building the global metrics.
	\item \textbf{Ensure that we capture as much data as possible.} In the case of wanting to develop the project on a larger scale, both providers can be contacted to try to obtain an increase in the consumption limit of their APIs. In this way, depending on the new limitation, it could try to divide its catch at a certain time. In the best case, it would be possible to capture all tweets, even if our system captures them periodically, since the identifier of the last tweet captured in the previous iteration can be used as a reference for the endpoint to start the new capture from that point.
	\item \textbf{Add visualizations in the front-end layer.} Additional visualizations could be introduced that would allow the user, for example, to group by artist and to obtain global metrics that answer questions such as: ``How popular are this artist's tracks?'' or ``How danceable are this artist's tracks compared to this other artist's tracks?''. Another possible view to develop would be one that allows tracks to be observed by metric ranges and answers questions such as what percentage of the total tracks are between 80-100 of the energy metric.
	\item \textbf{Upload from history functionality.} Another improvement that would facilitate user interaction with the tool would be the ability to load the file history from the tool itself. In this way, the user could stop the process, select his historical files and be able to load only the desired data.
	\item \textbf{Replace Docker Compose with a more suitable tool.} At the design level, although Docker Compose fits the requirements for this project, in the case that a large scalability is required, it would be convenient to replace the tool with another one that allows container management across multiple hosts, such as Docker Swarm or Kubernetes.
	\item \textbf{Refactor the code.} In the event that the process increases in the amount of data captured, it is advisable to perform an analysis of the capture process and a refactoring of the necessary parts to speed up the process, if necessary. With the current data volume, the execution speed of the whole process triggered by Apache Airflow takes an average of 30-35 seconds.
	\item \textbf{Partial or total migration to the cloud.} Finally, if deemed consistent with new requirements that may arise in the future, parts or even the project could be migrated to a cloud infrastructure. For example, the Data Warehouse could be replaced by one of the native offerings of the major Cloud providers.
\end{itemize}

\nonzeroparskip Most of these steps were considered for inclusion at some point in the development phase, but were discarded to ensure that the project, which for time reasons was conceived from the beginning as a proof-of-concept rather than an end-user oriented application, remained on schedule.


%\renewcommand\chaptername{Anexo}
%\renewcommand\thechapter{\Roman{chapter}}
%\setcounter{chapter}{0}

% Añadir entrada en el índice: Anexos
\appendix
\addcontentsline{toc}{part}{Appendixes}
\part*{Appendixes}

\apendice{Project Plan}

\section{Introduction}
\nonzeroparskip The project planning was decided in an initial meeting between the author and his tutor. It was based in an Agile methodology, with two-weeks \textit{sprints} and meetings between the author and his tutor conditioned to their availability.

\nonzeroparskip The project repository was stored in GitHub under the url \url{https://github.com/AdrianRiesco/Masters-Thesis-on-Big-Data}. Each \textit{sprint} was created as a \textit{milestone}, with the \textit{issues} contained there being the tasks assigned. The \textit{issues} were created to reflect tasks at most eight hours, allowing the author segregate his work and manage each \textit{sprint} better. The author closed an \textit{issue} when the task was finished and a \textit{milestone} when the \textit{sprint} was over, regardless of its state. If a task remained in an open state when a \textit{sprint} reached its planned end date, the \textit{issue} was transfered to the next \textit{milestone}.

\nonzeroparskip A meeting was held by the author and his tutor at the end of each sprint. During these meetings, both of them reviewed the state and development of the tasks of the corresponding sprint and planned the tasks of the next sprint. All the \textit{milestones} and \textit{issues} can be consulted in the project repository.

\section{Temporary planning}
\nonzeroparskip The sprints carried out for the development of the project are described below with their corresponding dates:
\begin{description}
	\item[Initial meeting.] Held on Monday January 31\textsuperscript{st}, it was the start point for the first sprint. During this meeting, the objective of the project, the data source and the tools to be used were validated by both the author and his tutor. The author previously made a research and came with an idea and the tutor exposed his point of view to create the final goal.
	\item[Sprint 1.] Weeks of January 31\textsuperscript{st} and February 7\textsuperscript{th}. This Sprint had the following tasks assigned:
	\begin{itemize}
		\item Configure the work environment.
		\item Configure the project memory template.
		\item Write a draft of the objectives and main goals.
		\item Write a brief description of the tools selected.
		\item Write a brief explanation of the selected tools and the work methodology.
		\item Inspect Twitter API.
		\item Inspect Spotify API.
	\end{itemize}
	The end-of-sprint meeting was held on Wednesday February 16\textsuperscript{th}.
	Analysis: Most of the activities were realized by the author, excepting the Inspection of the Spotify API. Regarding the Twitter API, the author inspected the output and he concluded that it had the characteristics needed to be used to launch queries to the Spotify API (the tweet could be cleaned to get the song name and artist).
	\item[Sprint 2.] Weeks of February 14\textsuperscript{th} and February 21\textsuperscript{st}. This Sprint had the following tasks assigned:
	\begin{itemize}
		\item Write the code to gather information from Spotify.
		\item Write the code to gather information from Twitter.
		\item Write a description of Spotify and Twitter APIs.
		\item Write the API inspection process in the ``Programmer guide'' section.
		\item Write the project introduction.
		\item Write the Twitter and Spotify data description.
	\end{itemize}
	The end-of-sprint meeting was held on Wednesday February 2nd. Analysis: All the activities were accomplished on time. The author identified some potential problems when extracting information from the Spotify API, such as the name of the song must be quite accurate to be able to get the search results.
	\item[Sprint 3.] Weeks of February 28\textsuperscript{th} and March 7\textsuperscript{th}. This Sprint had the following tasks assigned:
	\begin{itemize}		
		\item Write a description of Twitter API.
		\item Write a description of Spotify API.
		\item Write the description of the Twitter data.
		\item Write the description of the Spotify data.
		\item Improve and unify the code written to collect data from both APIs.
		\item Set up the Spark environment using Docker.
	\end{itemize}
	The end-of-sprint meeting was held on Wednesday March 16\textsuperscript{th}. Analysis: During this sprint, the author had difficulties configuring the Docker environment, which derived in a delay of the other tasks. These tasks were transfered to the next sprint.
	\item[Sprint 4.] Weeks of March 14\textsuperscript{th} and March 21\textsuperscript{st}. This Sprint had the following tasks assigned:
	\begin{itemize}
		\item Write a description of NoSQL Databases.
		\item Set up the Spark environment using Docker.
		\item Set up the Airflow environment using Docker.
	\end{itemize}
	The end-of-sprint meeting was held on Wednesday May 4\textsuperscript{th}. Analysis: Due to personal reasons, the author could not continue with the project during the month of April, so there was a temporary pause in the planning.
	\item[Sprint 5.] Weeks of May 2\textsuperscript{nd} and May 9\textsuperscript{th}. This Sprint had the following tasks assigned:
	\begin{itemize}
		\item Configure the DAG in Airflow.
		\item Learn the fundamentals of Flask.
		\item Redesign the project plan excluding the month of April
	\end{itemize}
	The end-of-sprint meeting was held on Wednesday May 18\textsuperscript{th}. Analysis: All tasks were accomplished on time.
	\item[Sprint 6.] Weeks of May 16\textsuperscript{th} and May 23\textsuperscript{rd}. This Sprint had the following tasks assigned:
	\begin{itemize}
		\item Deploy Cassandra environment.
		\item Write the description of the Twitter data.
		\item Write the description of the Spotify data.
		\item Write a description of an orchestrator.
		\item Write a description of CI/CD.
	\end{itemize}
	The end-of-sprint meeting was held on Monday May 30\textsuperscript{th}. Analysis: All tasks were accomplished on time. The Cassandra deployment took most of Sprint's time, generating a few prototypes before the final version.
	\item[Sprint 7.] Weeks of May 30\textsuperscript{th} and June 6\textsuperscript{nd}. This Sprint had the following tasks assigned:
	\begin{itemize}
		\item Integrate Cassandra with Airflow and Spark.
		\item Create the ETL workflow to load the data to Cassandra.
		\item Write a description of Apache Spark.
		\item Write a description of Docker and Docker Compose.
		\item Write a description of Flask, Jinja and Bootstrap.
	\end{itemize}
	The end-of-sprint meeting was held on Tuesday June 14\textsuperscript{th}. Analysis: All tasks were accomplished on time. The integration of Cassandra with Airflow and Spark took most of Sprint's time.
	\item[Sprint 8.] Weeks of June 13\textsuperscript{th} and June 20\textsuperscript{th}. This Sprint had the following tasks assigned:
	\begin{itemize}
		\item Create the front-end with Flask and Bootstrap.
		\item Integrate the ETL workflow with the front-end.
		\item Write section 5 ``Relevant aspects of the project''.
		\item Write section 7 ``Conclusions of the project''.
	\end{itemize}
	The end-of-sprint meeting was held on Monday June 27\textsuperscript{th}. An additional meeting was held in the middle of the sprint to review and propose changes to the first version of the front-end. Analysis: Although the front-end development was very time-consuming, all tasks were accomplished on time. At this point, all development work was finished.
	\item[Sprint 9.] Weeks of June 27\textsuperscript{rd} and July 4\textsuperscript{th}. This Sprint had the following tasks assigned:
	\begin{itemize}
		\item Write section 6 ``Related works''.
		\item Complete Appendix A ``Project plan''.
		\item Complete Appendix B ``Requirements''.
		\item Complete Appendix C ``Design''.
		\item Complete Appendix D ``Programmer's guide''.
		\item Complete Appendix E ``User manual''.
		\item Review the entire project report.
	\end{itemize}
	The end-of-sprint meeting was held on M--- July --\textsuperscript{th}. An additional mid-sprint meeting was held to review the report and plan the final tasks to finalize the project for defense.Analysis: The report was finished and the project was delivered.
\end{description}

\section{Feasibility study}
\nonzeroparskip The architecture of the project and the use case were designed to ensure its feasibility.

\subsection{Economic feasibility}
\nonzeroparskip The project is based on open-source platforms to ensure its economic and legal feasibility. The APIs where the information was gathered are free to use if the developer keeps his queries under specific limit rates.

\subsection{Legal feasibility}
\nonzeroparskip The project is based on open-source platforms to ensure its economic and legal feasibility. The licenses in which the technology and tools used in the project are based are:
\begin{itemize}
	\item \textbf{Twitter License.} Free with usage limitations and requires developer account~\cite{twitter_dev_license}.
	\item \textbf{Spotify License.} Free with usage limitations and requires developer account~\cite{spotify_dev_license}.
	\item \textbf{Apache License.} Free with limitations~\cite{apache_license}.
	\item \textbf{Flask License.} Free with limitations, uses BSD-3-Clause license~\cite{flask_license}.
	\item \textbf{Bootstrap License.} Free with limitations, uses Massachusetts Institute of Technology (MIT) license~\cite{bootstrap_license}.
	\item \textbf{\TeX{}maker License.} Free with limitations, uses GNU General Public License (GPL) license~\cite{texmaker}.
\end{itemize}

\tablaSmall{Tools and technologies' licenses}{l c c c c c c}{licensespertool}
{ \multicolumn{1}{l}{Tools} & Twitter & Spotify & Apache & BSD-3-Clause & MIT & GPL\\}{ 
Twitter API & X & & & & &\\
Spotify API & & X & & & &\\
Apache Airflow & & & X & & &\\
Apache Spark & & & X & & &\\
Apache Cassandra & & & X & & &\\
Flask & & & & X & &\\
Bootstrap & & & & & X &\\
Datatables & & & & & X &\\
Chart.js & & & & & X &\\
\TeX{}Maker & & & & & & X\\
}
\apendice{Requirements} \label{requirements}

\section{Introduction}

\nonzeroparskip This section lists the general objectives and requirements identified during the initial planning of the project and on whose fulfillment the development of the project has focused.

\section{General objectives}

\nonzeroparskip The requirements through which the use case was built were the following:
\begin{itemize}
	\item Ability to obtain data in real time.
	\item Combination of at least two different data sources.
	\item Potential to scale in both technology and data volume.
	\item Involvement of various technologies in the Big Data field.
	\item Use of open source tools.
\end{itemize}

\section{Catalog of requirements}

\subsection{Functional requirements}

\nonzeroparskip The functional requirements (FR) that the project had to meet were:
\begin{itemize}
	\item \textbf{FR1}. The data must be obtained from the Twitter hashtag \texttt{\textit{\#NowPlaying}} every 30 minutes, taking care of API rate limits.
	\begin{itemize}
		\item \textbf{Status}: Met.
		\item \textbf{Rationale}: The Apache Airflow DAG was configured with a run frequency of 30 minutes. The script used to collect the data has a maximum of 100 tweets collected per run to comply with API rate limits.
	\end{itemize}
	\item \textbf{FR2}. There must be at least two different visualizations and one of them must provide the ability to view all of the stored data.
	\begin{itemize}
		\item \textbf{Status}: Met.
		\item \textbf{Rationale}: There are two visualizations in the web application, ``Data'' and ``Visuals'', and the first one displays all the stored data.
	\end{itemize}
	\item \textbf{FR3}. At least one of the visualizations must show last songs name, artist and audio features.
	\begin{itemize}
		\item \textbf{Status}: Met.
		\item \textbf{Rationale}: The ``Data'' view of the web application displays all required fields.
	\end{itemize}
	\item \textbf{FR4}. At least one of the visualizations must have a link to the source tweet.
	\begin{itemize}
		\item \textbf{Status}: Met.
		\item \textbf{Rationale}: The ``Data'' view of the web application contains a link to the source tweet.
	\end{itemize}
	\item \textbf{FR5}. At least one of the visualizations must have the ability to compare different metrics.
	\begin{itemize}
		\item \textbf{Status}: Met.
		\item \textbf{Rationale}: The ``Visuals'' view of the web application provides the ability to visually compare metrics.
	\end{itemize}
	\item \textbf{FR6}. At least one of the visualizations must combine two different types of visualizations.
	\begin{itemize}
		\item \textbf{Status}: Met.
		\item \textbf{Rationale}: The ``Visuals'' view of the web application combines bar and line formats in the same chart.
	\end{itemize}
	\item \textbf{FR7}. Both visualizations must provide sorting capabilities.
	\begin{itemize}
		\item \textbf{Status}: Met.
		\item \textbf{Rationale}: Both ``Data'' and ``Visuals'' views of the web application provides provides the ability to sort by metric.
	\end{itemize}
	\item \textbf{FR8}. Both visualizations must be responsive to different screen sizes.
	\begin{itemize}
		\item \textbf{Status}: Met.
		\item \textbf{Rationale}: The web application was developed using Bootstrap as the CSS framework to facilitate scaling and resizing on different screens.
	\end{itemize}
\end{itemize}

\subsection{Technical requirements}

\nonzeroparskip The technical requirements (TR) that the project had to meet were:
\begin{itemize}
	\item \textbf{TR1}. The development must have the ability to be deployed in different environments with minimum effort.
	\begin{itemize}
		\item \textbf{Status}: Met.
		\item \textbf{Rationale}: The services were developed using multiple containers managed via Docker Compose.
	\end{itemize}
	\item \textbf{TR2}. The data flow must be automated, with the entire process orchestrated by a single tool.
	\begin{itemize}
		\item \textbf{Status}: Met.
		\item \textbf{Rationale}: All data flow is orchestrated by Apache Airflow.
	\end{itemize}
	\item \textbf{TR3}. The execution of the ETL process must be done with a tool that can scale and run in distributed environments.
	\begin{itemize}
		\item \textbf{Status}: Met.
		\item \textbf{Rationale}: Data extraction, transformation and loading (into a .csv file) is processed by Apache Spark, and data loading to Cassandra is performed by a Cassandra Query Language shell (cqlsh) command.
	\end{itemize}
	\item \textbf{TR4}. The data warehouse must have the ability to escalate in terms of a Big Data problem.
	\begin{itemize}
		\item \textbf{Status}: Met.
		\item \textbf{Rationale}: Apache Cassandra is the Data Warehouse used.
	\end{itemize}
	\item \textbf{TR5}. The web application must be designed with widely recognized tools.
	\begin{itemize}
		\item \textbf{Status}: Met.
		\item \textbf{Rationale}: The web application was designed with Flask and Bootstrap, and the visualizations with Chart.js and Datatables.
	\end{itemize}
	\item \textbf{TR6}. All the tools used must be open source.
	\begin{itemize}
		\item \textbf{Status}: Met.
		\item \textbf{Rationale}: All the tools used to develop the project are open source.
	\end{itemize}
\end{itemize}
\apendice{Design specification}

\section{Introduction}

\nonzeroparskip This section shows the structure of the data used and the flow and architecture diagrams.

\section{Data design}

\subsection{Twitter data structure}
\nonzeroparskip The data received from Twitter queries has the following structure (the data has been obtained from a real query and its output has been reduced by trimming certain elements due to their length):

\begin{lstlisting}[language=json,firstnumber=1]
{"data":[
      {
         "id":"1533311938209382403",
         "entities":{
            "annotations":[
               {
                  "start":43,
                  "end":62,
                  "probability":0.6061,
                  "type":"Other",
                  "normalized_text":"Breakfast In America"
               }
            ],
            "urls":[
               {
                  "start":84,
                  "end":107,
                  "url":"https://t.co/YiXxSepm8x",
                  "expanded_url":"https://rideshare.airtime.pro",
                  "display_url":"rideshare.airtime.pro",
                  "images":[
                     {
                        "url":"https://pbs.twimg.com/news_img/1532561096858640397/3mUiSDDN?format=jpg&name=orig",
                        "width":1920,
                        "height":1200
                     },
                     {
                        "url":"https://pbs.twimg.com/news_img/1532561096858640397/3mUiSDDN?format=jpg&name=150x150",
                        "width":150,
                        "height":150
                     }
                  ],
                  "status":200,
                  "title":"Rideshare Radio",
                  "description":"Hits from the 70's 80s 90s 00s 10s 20s, No Talking just back to back Music totally commercials free 24/7",
                  "unwound_url":"https://rideshare.airtime.pro"
               }
            ],
            "hashtags":[
               {
                  "start":0,
                  "end":11,
                  "tag":"Nowplaying"
               },
               {
                  "start":129,
                  "end":139,
                  "tag":"Rideshare"
               }
            ]
         },
         "created_at":"2022-06-05T04:57:08.000Z",
         "text":"#Nowplaying 2010 Remastered - Supertramp - Breakfast In America  - Stream here-&gt; https://t.co/YiXxSepm8x - Non Stop Hits 24/7 #Rideshare #Radio #Hits #Uber #RideshareRadio #Petrol #Parcoursup #PlatinumJubilee #ENGvNZ #PrideMonth"
      },
      {
         Second tweet.      
      }
   ],
   "meta":{
      "newest_id":"1533311938209382403",
      "oldest_id":"1533311921256124416",
      "result_count":10,
      "next_token":"b26v89c19zqg8o3fpyzltxkhapj3hf1q96mc01w4yl3el"
   }
}
\end{lstlisting}

\nonzeroparskip The fields required for the project are:
\begin{itemize}
	\item \textbf{\texttt{id.}} Tweet id, useful uniquely identify the tweet.
	\item \textbf{\texttt{text.}} Tweet text, useful to identify the song played.
	\item \textbf{\texttt{entities.}} useful to clean the text and remove hashtasg, cashtags, mentions and urls.
	\item \textbf{\texttt{created\_at.}} Tweet creation date.
\end{itemize}

\subsection{Spotify data structure}
\nonzeroparskip The data received from Spotify queries to Search endpoint (``Search for Item'') has the following structure (the data has been obtained from a real query and its output has been reduced by trimming certain elements due to their length):
\begin{lstlisting}[language=json,firstnumber=1]
{"tracks":{
      "href":"https://api.spotify.com/v1/search?query=Savoy+Brown+Wang+Dang+Doodle+&type=track&offset=0&limit=1",
      "items":[
         {
            "album":{
               "album_type":"album",
               "artists":[
                  {
                     "external_urls":{
                        "spotify":"https://open.spotify.com/artist/17obwOahRWI121iMUZznh2"
                     },
                     "href":"https://api.spotify.com/v1/artists/17obwOahRWI121iMUZznh2",
                     "id":"17obwOahRWI121iMUZznh2",
                     "name":"Savoy Brown",
                     "type":"artist",
                     "uri":"spotify:artist:17obwOahRWI121iMUZznh2"
                  }
               ],
               "available_markets":[
                  "MX",
                  "US"
               ],
               "external_urls":{
                  "spotify":"https://open.spotify.com/album/5oq20r8iNOO9fpw8R2h3vE"
               },
               "href":"https://api.spotify.com/v1/albums/5oq20r8iNOO9fpw8R2h3vE",
               "id":"5oq20r8iNOO9fpw8R2h3vE",
               "images":[
                  {
                     "height":640,
                     "url":"https://i.scdn.co/image/ab67616d0000b2735dd44bf0a252e30d4bb2e7c8",
                     "width":640
                  },
                  {
                     "height":300,
                     "url":"https://i.scdn.co/image/ab67616d00001e025dd44bf0a252e30d4bb2e7c8",
                     "width":300
                  }
               ],
               "name":"Street Corner Talking",
               "release_date":"1971-01-01",
               "release_date_precision":"day",
               "total_tracks":8,
               "type":"album",
               "uri":"spotify:album:5oq20r8iNOO9fpw8R2h3vE"
            },
            "artists":[
               {
                  "external_urls":{
                     "spotify":"https://open.spotify.com/artist/17obwOahRWI121iMUZznh2"
                  },
                  "href":"https://api.spotify.com/v1/artists/17obwOahRWI121iMUZznh2",
                  "id":"17obwOahRWI121iMUZznh2",
                  "name":"Savoy Brown",
                  "type":"artist",
                  "uri":"spotify:artist:17obwOahRWI121iMUZznh2"
               }
            ],
            "available_markets":[
               "MX",
               "US"
            ],
            "disc_number":1,
            "duration_ms":440733,
            "explicit":false,
            "external_ids":{
               "isrc":"GBF077120720"
            },
            "external_urls":{
               "spotify":"https://open.spotify.com/track/7p99XDR7dKaIMTYV3zia0V"
            },
            "href":"https://api.spotify.com/v1/tracks/7p99XDR7dKaIMTYV3zia0V",
            "id":"7p99XDR7dKaIMTYV3zia0V",
            "is_local":false,
            "name":"Wang Dang Doodle",
            "popularity":19,
            "preview_url":"None",
            "track_number":7,
            "type":"track",
            "uri":"spotify:track:7p99XDR7dKaIMTYV3zia0V"
         }
      ],
      "limit":1,
      "next":"https://api.spotify.com/v1/search?query=Savoy+Brown+Wang+Dang+Doodle+&type=track&offset=1&limit=1",
      "offset":0,
      "previous":"None",
      "total":11
   }
}
\end{lstlisting}

\nonzeroparskip The fields required for the project are:
\begin{itemize}
	\item \textbf{\texttt{id.}} Track id, useful uniquely identify the track.
	\item \textbf{\texttt{name.}} Track name, useful to identify the track.
	\item \textbf{\texttt{popularity.}} Popularity of the track.
	\item \textbf{\texttt{artists' id.}} ID of the artists.
	\item \textbf{\texttt{artists' name.}} Name of the artists.
\end{itemize}

\nonzeroparskip The data received from Spotify queries to the Tracks endpoint (``Get Tracks' Audio Features'') has the following structure (the data has been obtained from a real query and its output has been reduced by trimming certain elements due to their length):
\begin{lstlisting}[language=json,firstnumber=1]
{"audio_features":[
      {
         "danceability":0.516,
         "energy":0.36,
         "key":7,
         "loudness":-11.264,
         "mode":1,
         "speechiness":0.03,
         "acousticness":0.83,
         "instrumentalness":0.885,
         "liveness":0.116,
         "valence":0.144,
         "tempo":127.176,
         "type":"audio_features",
         "id":"7dg3XqARw7qOrkt9pZZNRF",
         "uri":"spotify:track:7dg3XqARw7qOrkt9pZZNRF",
         "track_href":"https://api.spotify.com/v1/tracks/7dg3XqARw7qOrkt9pZZNRF",
         "analysis_url":"https://api.spotify.com/v1/audio-analysis/7dg3XqARw7qOrkt9pZZNRF",
         "duration_ms":233812,
         "time_signature":4
      },
      {
         Group of features of the second track.
      }...]
}
\end{lstlisting}

\nonzeroparskip The fields required for the project are:
\begin{itemize}
	\item id.
	\item danceability.
	\item energy.
	\item key.
	\item loudness.
	\item mode.
	\item speechiness.
	\item acousticness.
	\item instrumentalness.
	\item liveness.
	\item valence.
	\item tempo.
	\item duration\_ms.
	\item time\_signature.
\end{itemize}

\subsection{Cleaned data}
\nonzeroparskip After the cleaning process, the resulting data structure is:

\begin{itemize}
	\item id\_tweet.
	\item text.
	\item created\_at.
	\item url\_tweet.
	\item id\_track.
	\item name.
	\item popularity.
	\item artists\_id.
	\item artists\_name.
	\item danceability
	\item energy
	\item key
	\item loudness
	\item mode
	\item speechiness
	\item acousticness
	\item instrumentalness
	\item liveness
	\item valence
	\item tempo
	\item duration\_ms
	\item time\_signature
\end{itemize}

\section{Architectural design}
\nonzeroparskip The component diagram is shown in the figure~\ref{component-diagram}. Each block represents a service deployed in one or more containers, except for data sources, which are external APIs. The UML format is referenced in the figure~\ref{component-diagram-flowchart}.

\figuraNormalSinMarco{0.22}{img/component-diagram}{Component diagram}{component-diagram}{}

\figuraNormalSinMarco{0.21}{img/component-diagram-flowchart}{UML component diagram}{component-diagram-flowchart}{}

\section{Procedural design}
\nonzeroparskip The flow diagram is shown in the figure~\ref{flow-diagram}.
\begin{enumerate}
	\item Get tweets with the hashtag \texttt{\#NowPlaying}.
	\item Search the track information and audio features.
	\item Get the track information and audio features.
	\item Load the data into the Data Warehouse.
	\item Request the data from the web application.
\end{enumerate}

\figuraNormalSinMarco{0.19}{img/flow-diagram}{Flow diagram}{flow-diagram}{}

\nonzeroparskip The sequence diagram is depicted in figure~\ref{sequence-diagram-flowchart}. The loop represents the ``spark\_main'' DAG, executed every 30 minutes and composed of three API calls, an Apache Cassandra load and the creation of the history file. The section outside the loop represents the user accessing the web application.

\figuraNormalSinMarco{0.46}{img/sequence-diagram-flowchart}{UML sequence diagram}{sequence-diagram-flowchart}{}
\apendice{Programming technical documentation}

\section{Introduction}

\section{Directory structure}
Estructura GitHub.

\section{Programmer's guide}

\subsection{Analysis}

\nonzeroparskip During the analysis phase, the author inspected the output of Twitter and Spotify APIs using Postman. In the first place, relying on the Twitter documentation, the author inspected the Twitter API by following the next steps:
\begin{enumerate}
	\item Get access to the Twitter Developer Portal.
	\item Get the credentials needed to consult the different endpoints of the API.
	\item Import the \textit{Twitter API v2} collection on Postman.
	\item Create a fork of the automatically created environment (\textit{Twitter API v2}) and collection \textit{Twitter API v2} to be able to edit the values.
	\item Modify the environment to include the following developer keys and tokens:
	\begin{itemize}
		\item Consumer key (\texttt{consumer\_key}).
		\item Consumer secret (\texttt{consumer\_secret}).
		\item Access token (\texttt{access\_token}).
		\item Token secret (\texttt{token\_secret}).
		\item Bearer token (\texttt{bearer\_token}).
	\end{itemize}
	\item In the collection tab, select the endpoint \textit{Search Tweets} $\longrightarrow$ \textit{Recent search} for the initial exploration. Configure the following parameters:
	\begin{itemize}
		\item \texttt{query} = \texttt{\#NowPlaying}
		\item \texttt{tweet.fields} = \texttt{created\_at,entities}
		\item \texttt{max\_results} = \texttt{10}
	\end{itemize}
	\item Now, we can send our query \url{https://api.twitter.com/2/tweets/search/re cent?query=\%23NowPlaying\&max\_results=10\&tweet.fields=created\_a t,entities} to get the 10 most recent tweets with the hashtag \#NowPlaying and receive their basic information (id, text) as well as the entities (hashtags, urls, annotations...) and the creation time stamp.
\end{enumerate}

\nonzeroparskip After analyze the data gathered from the Twitter API,  the author inspected the Spotify API, more specifically the endpoint ``Search for Item''), following the next steps:
\begin{enumerate}
	\item Get access to the Spotify Developer Portal.
	\item Enter in the developers console and select the ``Search for Item'' endpoint.
	\item Specify the parameters of the search. We can specify the type ``track'' and add a limit of one to only receive the first song found.
	\item After click on get the bearer token, we can use that token by clicking on Try Me or just copy the resulting query in our Linux console to check the output.
	\item The result of this query is the first result of the search containing information of the artist, the song and the album. With the artist and song ids we can consult other endpoints to get an audio analysis, the audio features and the artist information, between others.
\end{enumerate}

\subsection{Development}
\nonzeroparskip During the development phase, the following items were installed in the system:
\begin{itemize}
	\item Docker version 20.10.12, build e91ed57.
	\item docker-compose version 1.29.2, build 5becea4c.
	\item docker-compose version 1.29.2, build 5becea4c.
\end{itemize}

\nonzeroparskip Steps followed:
\begin{enumerate}
	\item Create the dockerfile.
	\item Build the image:  \texttt{docker build -f spark.Dockerfile -t ar/spark-ubuntu .}
	\item Create a Docker network: \texttt{docker network create --driver bridge spark-network}
	\item Run the image: \texttt{sudo docker run -d -t --name master --network spark-network ar/spark-ubuntu}
\end{enumerate}

\section{Compilation, installation and execution of the project}


\section{System tests}


\apendice{User documentation}

\section{Introduction}
\nonzeroparskip This section summarizes the elements that the user must have and the steps that must be followed to correctly execute the project.

\section{User requirements}
\nonzeroparskip The technical user requirements to run the project are defined in the section~\ref{programmer_development}. As a summary:
\begin{enumerate}
	\item Docker.
	\item Docker Compose.
	\item Python 3.
\end{enumerate}

\nonzeroparskip In addition to the technical requirements, the user must have at least a basic/intermediate knowledge level of the operating system that will be used to run the project.

\fbox{\parbox{0.9\textwidth}{\textit{The commands described for the installation and usage of the project are oriented to Linux environments.}}}

\section{Installation}
\nonzeroparskip The steps that need to be followed to install the project are referenced in the section~\ref{programmer_execute}. A simpler summary is shown below:
\begin{enumerate}
	\item Open a \texttt{command prompt} and execute the following commands to set up the permissions of Airflow required folders:
	\begin{itemize}
		\item \texttt{sudo chmod -R 777 spark/resources/history}
		\item \texttt{sudo chmod -R 777 airflow/logs}
		\item \texttt{sudo chmod -R 777 airflow/plugins}
	\end{itemize}
	\item Move to the \texttt{docker} folder and launch the environment with ``\texttt{sudo docker-compose up --build -d}''.
	\item Wait until all the services have started and are in a healthy state.
	\item Access to the Airflow UI (\url{http://localhost:8080}, user ``\texttt{airflow}'' and password ``\texttt{airflow}'') and start the DAG ``\texttt{spark\_main}'' as shown in the figure~\ref{start-dag}.
	\item Use the command ``\texttt{sudo docker-compose down}'' to stop the process.
\end{enumerate}

\section{User's manual} \label{usermanual}

\nonzeroparskip The usage of the project is simple. After installing and running the project, including the activation of the Airflow DAG, all data results are presented to the user in the web application (\url{http://localhost:8000}), containing three views.
\begin{itemize}
	 \item \textbf{Home.} It shows a brief introduction of the project and an animated gif to illustrate the implemented data flow, along with a link to the repository on GitHub. The ``Home'' view and its components are referenced in figure \ref{front-home-described}.
	 \item \textbf{Data.} It displays a table with all the data extracted. The table has been made with Datatables and allows the user to:
	 \begin{itemize}
	 	\item \textbf{Column selector}. Allows to select the columns to be hidden or shown. The default view is shown with all columns hidden except ``Popularity'', ``Danceability'', ``Energy'' and ``Duration(s)''. The hidden columns are grayed out.
	 	\item \textbf{Size selector}. Allows to select the amount of data to be displayed, limited to 10 (default), 25, 50 and 100.
	 	\item \textbf{Search bar}. Allows to search the entire dataset, regardless of page or hidden columns.
	 	\item \textbf{Page selector}. Allows to select the page to be displayed (first by default).
	 	\item \textbf{Number of entries}. Shows the range of rows displayed on the page and the total number of entries in the dataset.
	 	\item \textbf{Data displayed}. The displayed data can be sorted by clicking on the column headers. The default sorting is by date from most recent to least recent.
	 \end{itemize}
	 The ``Data'' view and its components are referenced in figure \ref{front-data-described}.
	 \item \textbf{Visuals.} It displays a chart composed by 
	 with two data selectors, a multi-select to add data as a bar and a single-select to add data as a line. A script was added that does not allow both selectors to have the same column in their data, so when a column that is present in one of them is selected in the other, it is removed from the previous one. It also has a selector to choose the amount of data to display (values of 5, 10, 15 and 20, to avoid saturating the graph), as well as an order selector that allows the user to sort by any of the columns in ascending or descending order.
	 \begin{itemize}
	 	\item \textbf{Quantity selector}. Allows to select the amount of data to be displayed, limited to 5, 10 (default), 15 and 20.
	 	\item \textbf{Order selector}. Allows to select the data sorting, with a drop-down to select the column and a switch to select the direction (ascending or descending). The default sorting is by date from most recent to least recent.
	 	\item \textbf{Column selector for the bars}. Allows to select the columns to hide or show in the bar format. The default view displays only the ``Popularity'' column.
	 	\item \textbf{Column selector for the line}. Allows to select the columns to hide or show in the line format. The default view does not display any columns (selected).
	 	\item \textbf{Data displayed}. The displayed data can be hidden by clicking on the labels shown at the top.
	 \end{itemize}
	 The ``Visuals'' view and its components are referenced in figure \ref{front-visuals-described}.
\end{itemize}

\figuraNormalSinMarco{0.25}{img/front-home-described}{Front-end home view}{front-home-described}{}

\figuraNormalSinMarco{0.25}{img/front-data-described}{Front-end data view}{front-data-described}{}

\figuraNormalSinMarco{0.23}{img/front-visuals-described}{Front-end graph view}{front-visuals-described}{}


\bibliographystyle{plain}
\bibliography{bibliography}

\end{document}
